\documentclass[11pt,a4paper]{article}
\usepackage[T1]{fontenc}
\usepackage[utf8]{inputenc}
\usepackage{lmodern}
\usepackage{amsmath,amssymb,amsthm,mathtools,mathrsfs}
\usepackage{graphicx,float,xcolor,multicol}
\usepackage[shortlabels]{enumitem}
\usepackage[margin=27mm]{geometry}
\usepackage{microtype,needspace,placeins}
\usepackage{tikz,tikz-cd}
\usetikzlibrary{arrows.meta,calc,positioning,patterns,decorations.pathreplacing}
\usepackage[colorlinks=true,linkcolor=teal,urlcolor=teal,citecolor=teal]{hyperref}
\setlength{\parindent}{0pt}
\setlength{\parskip}{.6em}
\setcounter{secnumdepth}{0}
\allowdisplaybreaks
\raggedbottom
\providecommand{\cross}{\times}
\DeclareUnicodeCharacter{2020}{\ensuremath{\dagger}}
\DeclareMathOperator{\supp}{supp}
\DeclareMathOperator*{\esssup}{ess\,sup}
\providecommand{\chapter}{\section}
\newenvironment{tool}[2]{\subsection*{Tool #1 --- #2}}{}
\newcommand{\ToolRef}[1]{Tool #1}
\newcommand{\FigureTag}[2]{\textit{Figure #1}}
\newcommand{\Result}[1]{\subsection*{#1}}
\newcommand{\Application}[1]{\subsubsection*{Application\if\relax\detokenize{#1}\relax\else\space--- #1\fi}}
\newcommand{\ProofHeading}{\par\textit{Proof.}\space}
\newcommand{\ProofOf}[1]{\par\textit{Proof of #1.}\space}
\newcommand{\RemarkHeading}{\par\textbf{Remark.}\space}
\newcommand{\NoteHeading}{\par\textbf{Note.}\space}
\newcommand{\ExerciseHeading}{\par\textbf{Exercise.}\space}

% Rendering support only; mathematical text is in each independent post.
\usepackage{booktabs,array,longtable}
\definecolor{HandoutBlue}{HTML}{2F6690}
\definecolor{HandoutNavy}{HTML}{17365D}
\newtheorem{theorem}{Theorem}
\newtheorem{lemma}[theorem]{Lemma}
\newtheorem{proposition}[theorem]{Proposition}
\newtheorem{corollary}[theorem]{Corollary}
\newtheorem{conjecture}[theorem]{Conjecture}
\newtheorem{definition}{Definition}
\newtheorem{example}{Example}
\newtheorem{namedtoolcounter}{Tool}
\newenvironment{namedtool}[1]{\begin{namedtoolcounter}[#1]}{\end{namedtoolcounter}}
\newtheorem*{claim}{Claim}
\newtheorem*{remark}{Remark}
\newtheorem*{exercise}{Exercise}
\newtheorem*{question}{Question}
\newtheorem*{observation}{Observation}
\newcommand{\SourceChapter}[2]{\section*{#2}}
\newcommand{\Lecture}[2]{\section*{Lecture #1: #2}}
\newcommand{\unclear}[1]{\textit{[unclear: #1]}}
\newcommand{\sourcegap}{\textit{[unfinished in the source]}}
\DeclareMathOperator{\dist}{dist}
\DeclareMathOperator{\id}{id}
\DeclareMathOperator{\Hol}{Hol}
\DeclareMathOperator{\Per}{Per}
\DeclareMathOperator{\Fix}{Fix}
\DeclareMathOperator{\diam}{diam}
\DeclareMathOperator{\Var}{Var}
\DeclareMathOperator{\Leb}{Leb}
\DeclareMathOperator{\Hom}{Hom}
\DeclareMathOperator{\End}{End}
\DeclareMathOperator{\Ann}{Ann}
\DeclareMathOperator{\Spec}{Spec}
\DeclareMathOperator{\Max}{Max}
\DeclareMathOperator{\Ob}{Ob}
\DeclareMathOperator{\Tor}{Tor}
\DeclareMathOperator{\Ass}{Ass}
\DeclareMathOperator{\Mor}{Mor}
\DeclareMathOperator{\Fun}{Fun}
\DeclareMathOperator{\Nat}{Nat}
\DeclareMathOperator{\Cone}{Cone}
\DeclareMathOperator{\MCG}{MCG}
\DeclareMathOperator{\Aut}{Aut}
\DeclareMathOperator{\Deck}{Deck}
\DeclareMathOperator{\SL}{SL}
\DeclareMathOperator{\PSL}{PSL}
\DeclareMathOperator{\Area}{Area}
\DeclareMathOperator{\im}{im}
\DeclareMathOperator{\PSh}{PSh}
\newcommand{\CRing}{\mathsf{CRings}}
\newcommand{\AMod}{A\text{-}\mathsf{Mod}}
\newcommand{\Ab}{\mathsf{Ab}}
\newcommand{\Z}{\mathbb Z}
\newcommand{\Q}{\mathbb Q}
\newcommand{\R}{\mathbb R}
\newcommand{\C}{\mathbb C}
\newcommand{\N}{\mathbb N}
\newcommand{\kk}{\mathbb k}
\renewcommand{\aa}{\mathfrak a}
\newcommand{\bb}{\mathfrak b}
\newcommand{\pp}{\mathfrak p}
\newcommand{\qq}{\mathfrak q}
\newcommand{\mm}{\mathfrak m}
\newcommand{\nn}{\mathfrak n}
\newcommand{\rad}{\mathrm r}
\newcommand{\cat}[1]{\mathsf{#1}}
\setcounter{secnumdepth}{0}
\allowdisplaybreaks

\graphicspath{{figures/ergodic-theory/}}
\begin{document}
\section*{Mixing, Products, and Spectral Criteria}
% BLOG-CONTENT-BEGIN
\subsection{Mixing}
Ergodic theorems give you all the equidistribution theorems one would hope
for, so let us now abandon the study of orbits. Recall that we studied
ergodicity (the contents of Lecture 2):
\begin{itemize}
 \item Measure-theoretic: to obtain dynamical information.
 \item Operator-theoretic: to obtain equidistribution theorems.
\end{itemize}
Now let us focus on the study of ``set-orbits''
$\{T^{\circ n}(E)\}_{n\in\mathbb Z}$. The measure-theoretic perspective
gave us information about $\{T^{\circ n}(x)\}_{n\in\mathbb Z}$, namely
recurrence theorems; the operator-theoretic perspective told us about
$\{f\circ T^n\}_{n\in\mathbb Z}$. Now it is time to give up these
perspectives and look at it for what it is: a probability space.
Let us view the system $(X,\mathcal B,\mu,T)$ probabilistically.
Before we jump into that, here is a small observation. We saw that
\[
 \frac1N\sum_{n=0}^{N-1}f\circ T^n\longrightarrow\int f\,d\mu
 \quad\text{in $L^2_\mu$, as $N\to\infty$.}
\]
One can restate this slightly to obtain a probabilistic interpretation
% Source page 2
of ergodicity:
\[
 T\text{ is ergodic}
 \quad\Longleftrightarrow\quad
 \frac1N\sum_{n=0}^{N-1}\mu(A\cap T^{-n}(B))
 \longrightarrow\mu(A)\mu(B).
\]
\begin{itemize}
 \item If $T$ is ergodic, then, by the dominated convergence theorem,
 \[
  \frac1N\sum_{n=0}^{N-1}\langle f\circ T^n,g\rangle
  \longrightarrow\left(\int f\,d\mu\right)\left(\int g\,d\mu\right).
 \]
 Set $f=\chi_B$ and $g=\chi_A$ to obtain the desired result.
 \item Conversely, if $T^{-1}(A)=A$, set $B=X\setminus A$. We obtain
 $\mu(A)\mu(X\setminus A)=0$, so $\mu(A)\in\{0,1\}$.
\end{itemize}
Thus, on average, $T$ pulls back an event and \emph{tries} to make it
independent of $A$. If $T^{-n}(B)$ and $A$ were independent, then
\[
 \mu(A\cap T^{-n}(B))
 =\mu(A)\mu(T^{-n}(B))=\mu(A)\mu(B),
\]
and convergence would be assured. In general, however, an ergodic map just
tries to separate events and is not always successful. Keeping this in
mind, let us look at maps that asymptotically make $T^{-n}(B)$ and $A$
independent.

\setcounter{definition}{5}\begin{definition}[Mixing]
A system $(X,\mathcal B,\mu,T)$ is \emph{mixing} if, for all
$A,B\in\mathcal B$,
\[
 \mu(A\cap T^{-n}(B))\longrightarrow\mu(A)\mu(B)
 =\mu(A)\mu(T^{-n}(B))\qquad(n\to\infty).
\]
\end{definition}
\setcounter{definition}{6}\begin{definition}[Weak mixing]
A system $(X,\mathcal B,\mu,T)$ is \emph{weakly mixing} if, for all
$A,B\in\mathcal B$,
\[
 \frac1N\sum_{n=0}^{N-1}
 \bigl|\mu(A\cap T^{-n}(B))-\mu(A)\mu(B)\bigr|
 \longrightarrow0\qquad(N\to\infty).
\]
\end{definition}

% Source page 3
\begin{remark}
For a sequence $\{a_n\}_n$ (for example, $a_n$ may be
$\mu(A\cap T^{-n}B)$),
\[
 \lim_{n\to\infty}a_n=0
 \quad\Longrightarrow\quad
 \lim_{n\to\infty}\frac1n\sum_{i=0}^{n-1}|a_i|=0,
\]
but not conversely. Thus mixing implies weak mixing. Further, given a
weakly mixing system $(X,\mathcal B,\mu,T)$ and $T^{-1}B=B$, we have
\[
 \frac1N\sum_{n=0}^{N-1}|\mu(A\cap B)-\mu(A)\mu(B)|\longrightarrow0.
\]
Set $A=X\setminus B$. Then $\mu(A)\mu(B)\to0$, so $\mu(A)=0$ or
$\mu(B)=0$, and $T$ is ergodic. Consequently,
\[
 \text{Mixing}\ \subset\ \text{Weak mixing}\ \subset\ \text{Ergodicity}.
\]
Said differently:
\begin{itemize}
 \item $T^{-1}$ makes $A$ and $B$ asymptotically independent:
 this is mixing.
 \item $T^{-1}$ makes $A$ and $B$ asymptotically independent except on
 a zero-density set: this is weak mixing.
 \item $T^{-1}$ \emph{tries} to make $A$ and $B$ independent on average:
 this is ergodicity.
\end{itemize}
\end{remark}
\begin{remark}
Since $T^{-1}$ tries to separate events in whatever manner we prescribe,
we say that the inverse of $T^{-1}$, namely $T$, mixes events. Hence we
call such systems mixing or weakly mixing. It turns out that weak mixing
(a formulation due to von Neumann) is quite a natural thing, as it has so
many equivalent formulations. A generalisation of this concept will be
crucial for us in establishing Szemer\'edi's theorem.
\end{remark}
Now let us see other equivalent formulations of weak mixing, to find
% Source page 4
connections between the probabilistic interpretation and the others.

\setcounter{theorem}{14}\begin{theorem}[Mixing equivalences]
\label{thm:l4-mixing-equivalences}
Given a system $(X,\mathcal B,\mu,T)$, the following are equivalent:
\begin{enumerate}[label=(\arabic*)]
 \setcounter{enumi}{0}\item $T$ is weakly mixing.
 \setcounter{enumi}{1}\item $T\times T$ is ergodic with respect to $\mu\times\mu$.
 \setcounter{enumi}{2}\item $T\times T$ is weakly mixing with respect to $\mu\times\mu$.
 \setcounter{enumi}{3}\item If $(Y,\mathcal B',\nu,S)$ is ergodic, then
 $(X\times Y,\mathcal B\times\mathcal B',\mu\times\nu,T\times S)$
 is ergodic.
 \setcounter{enumi}{4}\item $U_T$ has \emph{continuous spectrum}: its only eigenvalue is $1$,
 and the corresponding eigenspace consists of constant functions.
\end{enumerate}
\end{theorem}

\setcounter{theorem}{15}\begin{lemma}[Koopman--Neumann]
\label{lem:l4-koopman-neumann}
Let $\{a_n\}$ be a sequence of bounded nonnegative real numbers.
The following are equivalent:
\begin{enumerate}[label=(\roman*)]
 \setcounter{enumi}{0}\item $\displaystyle\lim_{n\to\infty}\frac1n\sum_{j=0}^{n-1}a_j=0$.
 \setcounter{enumi}{1}\item There exists $J\subseteq\mathbb N$ of density zero, where
 \[
  \operatorname{density}(J):=
  \lim_{n\to\infty}\frac1n|J\cap[1,n]|,
 \]
 such that $a_n\to0$ as $n\to\infty$ with $n\notin J$.
 \setcounter{enumi}{2}\item $\displaystyle\lim_{n\to\infty}\frac1n\sum_{j=0}^{n-1}a_j^2=0$.
\end{enumerate}
\end{lemma}

% Source page 5


Let us stop this session by setting up the premise for the infamous
Rokhlin problem.
\setcounter{definition}{7}\begin{definition}[$k$-fold mixing]
A system $(X,\mathcal B,\mu,T)$ is \emph{$k$-fold mixing} if
\[
 \mu(A_0\cap T^{-n_1}A_1\cap\cdots\cap T^{-n_k}A_k)
 \longrightarrow\mu(A_0)\cdots\mu(A_k)
\]
as $n_1,n_2-n_1,n_3-n_2,\ldots,n_k-n_{k-1}\to\infty$, for all
$A_i\in\mathcal B$, $0\leq i\leq k$.
\end{definition}
\begin{question}[Rokhlin problem (open)]
If $(X,\mathcal B,\mu,T)$ is mixing, is it $k$-fold mixing for every
$k\geq1$?
\end{question}
% The source calls this problem open; no attempt has been made to update its mathematical status.
% BLOG-CONTENT-END
\end{document}
